\newpage
\section*{Solucionario completo}

\noindent Cada respuesta incluye la clave y una justificación breve. En los ejercicios numéricos se indica la ecuación o el cálculo principal.

\subsection*{Bloque A. Fundamentos teóricos y conceptuales}

\sol{1}{B}{La función de producción describe la frontera tecnológica: la cantidad máxima de producto obtenible con combinaciones dadas de insumos. No es una función de costos ni de beneficios.}

\sol{2}{C}{Una isocuanta agrupa todas las combinaciones de $K$ y $L$ que generan el mismo nivel de producción. La isocosto, en cambio, agrupa combinaciones de igual costo.}

\sol{3}{A}{La RMST indica cuánto capital puede reducirse al aumentar una unidad de trabajo sin alterar la producción. Es la pendiente (en valor absoluto) de la isocuanta.}

\sol{4}{C}{$RMST_{LK}=PMg_L/PMg_K$. En el óptimo de costos se iguala a $w/r$.}

\sol{5}{C}{Una RMST decreciente significa que cada vez se necesita menos capital para compensar una unidad adicional de trabajo; las isocuantas son convexas al origen.}

\sol{6}{D}{Rendimientos constantes a escala: $f(\lambda K,\lambda L)=\lambda f(K,L)$. Al duplicar los insumos, el producto se duplica.}

\sol{7}{B}{En $Q=AK^\alpha L^\beta$, al multiplicar los insumos por $\lambda$ se obtiene $\lambda^{\alpha+\beta}Q$. El grado de homogeneidad es $\alpha+\beta$.}

\sol{8}{B}{Si $\alpha+\beta>1$, el producto crece más que proporcionalmente al aumentar todos los insumos: rendimientos crecientes a escala.}

\sol{9}{D}{La ley de rendimientos marginales decrecientes es un fenómeno de corto plazo: con al menos un factor fijo, el producto marginal del factor variable eventualmente disminuye.}

\sol{10}{C}{El producto medio alcanza su máximo cuando se iguala al producto marginal. Si $PMg>PMe$, el $PMe$ aumenta; si $PMg<PMe$, disminuye.}

\sol{11}{B}{La etapa II comienza en el máximo del producto medio y termina cuando el producto marginal se hace cero. Es la región económicamente relevante.}

\sol{12}{B}{De $C=wL+rK$ se obtiene $K=C/r-(w/r)L$; la pendiente es $-w/r$.}

\sol{13}{A}{En el óptimo interior, la isocuanta es tangente a la isocosto: $RMST_{LK}=w/r$.}

\sol{14}{A}{La senda de expansión une las combinaciones óptimas de insumos para distintos niveles de producción, manteniendo los precios de los insumos constantes.}

\sol{15}{B}{La demanda condicionada de factores se deriva de minimizar $wL+rK$ sujeto a $Q=f(K,L)$. Depende de $Q$, $w$ y $r$.}

\sol{16}{B}{El lema de Shephard para costos: $\partial C(w,r,Q)/\partial w=L(w,r,Q)$. La derivada del costo respecto del precio del insumo es la demanda condicionada.}

\sol{17}{C}{Economías de escala: el costo medio disminuye al aumentar la producción. No implica que el costo total disminuya.}

\sol{18}{D}{Economías de alcance: $C(Q_1,Q_2)<C(Q_1,0)+C(0,Q_2)$. Producir conjuntamente es más barato que hacerlo por separado.}

\sol{19}{C}{Un costo hundido ya se incurrió y no es recuperable. No debe influir en decisiones marginales futuras, aunque puede afectar la decisión inicial de entrar.}

\sol{20}{C}{El costo económico incluye costos explícitos (pagos efectivos) y costos implícitos o de oportunidad, como el uso de capital propio o el tiempo del empresario.}

\sol{21}{A}{El beneficio normal es el rendimiento mínimo exigido por el capital propio; forma parte del costo de oportunidad y, por ello, del costo económico.}

\sol{22}{A}{En el corto plazo, si $P<CVMe$, la empresa no cubre ni los costos variables y le conviene suspender. Si $CVMe<P<CMe$, cubre variables y parte de los fijos, por lo que sigue operando.}

\sol{23}{B}{En el largo plazo no hay costos fijos irreversibles. Si $P<CMe$, la empresa obtiene pérdidas y sale del mercado.}

\sol{24}{C}{La empresa competitiva enfrenta $P=IMg$; maximiza beneficio donde $IMg=CMg$, es decir, $P=CMg$.}

\sol{25}{B}{La curva de oferta de corto plazo es el tramo de $CMg$ por encima del mínimo $CVMe$. Por debajo de ese precio, la empresa suspende.}

\sol{26}{A}{En el equilibrio de largo plazo, la libre entrada y salida lleva a $P=CMg=CMe$ y beneficio económico nulo; las empresas obtienen solo beneficio normal.}

\sol{27}{D}{El excedente del productor es el ingreso total menos el costo variable. Equivale al beneficio económico más el costo fijo.}

\sol{28}{D}{La dualidad establece que la función de costos contiene la misma información tecnológica que la función de producción, bajo condiciones de regularidad.}

\sol{29}{A}{Si $Q$ es homogénea de grado uno, Euler implica $Q=PMg_K K+PMg_L L$: el producto se agota si cada factor recibe su producto marginal.}

\sol{30}{B}{La elasticidad de sustitución mide la respuesta porcentual de la relación $K/L$ ante un cambio porcentual de la RMST: $\sigma=\dfrac{\%\Delta(K/L)}{\%\Delta RMST}$.}

\subsection*{Bloque B. Aplicaciones cuantitativas y optimización}

\sol{31}{C}{$Q=16^{0.5}25^{0.5}=4\cdot5=20$.}

\sol{32}{A}{$Q=2(9^{0.5})(4^{0.5})=2(3)(2)=12$.}

\sol{33}{C}{En el óptimo, $RMST=PMg_L/PMg_K=K/L=w/r=4$, de donde $K=4L$. Con $Q=20$: $Q=\sqrt{4L^2}=2L=20$, entonces $L=10$ y $K=40$. El costo es $4(10)+1(40)=80$.}

\sol{34}{A}{$K/L=w/r=4/1=4$.}

\sol{35}{D}{La suma de exponentes es $0.3+0.6=0.9<1$: rendimientos decrecientes a escala.}

\sol{36}{D}{La suma de exponentes es $0.4+0.7=1.1>1$: rendimientos crecientes a escala.}

\sol{37}{C}{En Leontief, $2K=L=40$. Entonces $K=20$ y $L=40$. El costo es $2(40)+3(20)=80+60=140$.}

\sol{38}{D}{Se requiere $K=30$ y $3L=30$, es decir $L=10$. El costo es $1(10)+2(30)=10+60=70$.}

\sol{39}{C}{El capital es más barato por unidad ($r=1<w=2$). Se usa solo capital: $K=50$, $L=0$; costo $=50$.}

\sol{40}{C}{$PMg_K/r=2/2=1$ y $PMg_L/w=1/3$. Conviene usar solo capital: $2K=40$, entonces $K=20$ y el costo es $2(20)=40$.}

\sol{41}{A}{$PMg_L=20L-3L^2$. En $L=4$: $80-48=32$.}

\sol{42}{B}{$Q=10(16)-64=96$; $PMe=Q/L=96/4=24$.}

\sol{43}{C}{$PMe=Q/L=10L-L^2$. Su máximo se obtiene con $dPMe/dL=10-2L=0$, es decir $L=5$.}

\sol{44}{D}{$CMg=dC/dQ=4+2Q$. En $Q=10$: $4+20=24$.}

\sol{45}{D}{$CMe=C/Q=100/Q+4+Q$. En $Q=10$: $10+4+10=24$.}

\sol{46}{B}{$CV=4Q+Q^2$. En $Q=10$: $40+100=140$; $CVMe=140/10=14$.}

\sol{47}{D}{$CMe=100/Q+4+Q$. Derivando: $-100/Q^2+1=0$, de donde $Q^2=100$ y $Q=10$.}

\sol{48}{A}{$P=CMg$: $20=4+2Q$, entonces $Q=8$.}

\sol{49}{B}{$IT=20(8)=160$. $CT=100+4(8)+8^2=100+32+64=196$. Beneficio $=160-196=-36$.}

\sol{50}{C}{$CVMe=4+Q$, que se minimiza en $Q=0$ con valor 4. Si $P<4$, la empresa no cubre ni los costos variables y suspende.}

\sol{51}{B}{En largo plazo, $P=\min CMe=24$ y cada empresa produce $Q=10$. La demanda de mercado es $Q=500-10(24)=260$. El número de empresas es $260/10=26$.}

\sol{52}{B}{$PS=\int_0^8[20-(4+2Q)]dQ=16(8)-8^2=128-64=64$.}

\sol{53}{A}{La función es CES con $\rho=0.5$. Entonces $\sigma=1/(1-\rho)=1/0.5=2$.}

\sol{54}{D}{Aquí $\rho=-1$; $\sigma=1/(1-(-1))=0.5$.}

\sol{55}{A}{Para $Q=K^{0.5}L^{0.5}$, $C=2Q\sqrt{wr}$. Con $w=4$ y $r=1$: $C=2Q(2)=4Q$.}

\sol{56}{B}{Con $K=16$, $Q=4\sqrt{L}$. Para $Q=8$: $\sqrt{L}=2$, $L=4$. El costo es $4(4)+1(16)=32$.}

\sol{57}{B}{De $C=Q^2/4+16$, $CMg=Q/2$. En $Q=8$: $CMg=4$.}

\sol{58}{D}{Si $C=4Q$, el costo marginal es constante e igual a 4.}

\sol{59}{C}{El costo unitario satisface $C(2Q)/C(Q)=2^{-0.2}=0.8706$. Cae aproximadamente $1-0.8706=0.1294$, es decir, 13\%.}

\sol{60}{D}{En $Q=K^{0.5}L^{0.5}$, la elasticidad del producto respecto de $L$ es el exponente 0.5.}

\sol{61}{D}{Con rendimientos crecientes a escala, el costo total crece menos que proporcionalmente que la producción: la elasticidad costo-producto es menor que 1.}

\sol{62}{B}{$L=Q\sqrt{r/w}$. La elasticidad respecto de $r$ es $0.5$: un aumento de 1\% en $r$ eleva $L$ en 0.5\%.}

\sol{63}{C}{Si $A$ aumenta 10\% y $K,L$ no cambian, $Q$ aumenta exactamente 10\%, porque $Q$ es proporcional a $A$.}

\sol{64}{A}{$K/L=w/r=1/4$, entonces $K=0.25L$. $Q=\sqrt{0.25L^2}=0.5L=10$, de donde $L=20$ y $K=5$. El costo es $1(20)+4(5)=40$.}

\sol{65}{B}{Con $K=16$, $Q=4\sqrt{L}=20$, entonces $\sqrt{L}=5$ y $L=25$. El costo es $4(25)+1(16)=116$.}

\subsection*{Bloque C. Casos con implicancias tarifarias y regulatorias}

\sol{66}{B}{$C(50)=1000+2(50)+0.1(50)^2=1000+100+250=1350$. $CMe=1350/50=27$.}

\sol{67}{D}{$CMg=dC/dQ=2+0.2Q$. En $Q=50$: $2+10=12$.}

\sol{68}{C}{$CMe=1000/Q+2+0.1Q$. Derivando: $-1000/Q^2+0.1=0$, de donde $Q^2=10000$ y $Q=100$.}

\sol{69}{C}{Con $P=CMg=12$ y $Q=50$, el ingreso es $12(50)=600$ y el costo es $1350$. La pérdida es $1350-600=750$.}

\sol{70}{A}{Un monopolio natural se caracteriza por costos medios decrecientes en el rango relevante y una función de costos subaditiva: una sola empresa abastece el mercado a menor costo que varias.}

\sol{71}{C}{Las economías de alcance surgen cuando producir agua y alcantarillado conjuntamente es más barato que producir cada servicio por separado.}

\sol{72}{A}{En RPI-X, la tarifa se ajusta por inflación menos un factor de eficiencia. Si la empresa reduce costos por debajo de lo previsto, retiene la diferencia hasta el siguiente ajuste.}

\sol{73}{A}{El efecto Averch-Johnson sostiene que, al remunerar el capital, la empresa regulada tiende a sobrecapitalizar para elevar la base sobre la que obtiene retorno.}

\sol{74}{D}{En cost-plus, el regulador reconoce los costos incurridos; la empresa tiene pocos incentivos a reducirlos y puede inflarlos.}

\sol{75}{B}{El yardstick competition compara la empresa con un grupo de referencia o con empresas similares, introduciendo presión competitiva sin duplicar la red.}

\sol{76}{A}{Si la energía se encarece, en el largo plazo la empresa sustituye energía por capital (equipos más eficientes) para reducir el costo por m$^3$.}

\sol{77}{D}{Productividad media del trabajo $=Q/L=1,000,000/50=20,000$ m$^3$ por trabajador.}

\sol{78}{C}{Crecimiento de la PTF $\approx$ crecimiento del producto menos crecimiento de los insumos $=6\%-4\%=2\%$.}

\sol{79}{D}{Si 25\% del agua producida no se factura, se factura 75\%. El costo por m$^3$ facturado es $1.50/0.75=2.00$.}

\sol{80}{B}{Con ANF de 15\%, el costo por m$^3$ facturado es $1.50/0.85=1.7647$. La reducción es $2.00-1.7647=0.2353\approx0.24$.}

\sol{81}{C}{El ahorro valorizado es $50,000\times2=100,000$. El beneficio neto es $100,000-80,000=20,000$.}

\sol{82}{B}{En hora punta se requiere capacidad adicional y el costo marginal de suministro es mayor; la tarifa de punta debe reflejarlo para incentivar el uso eficiente.}

\sol{83}{A}{La inversión hundida crea un problema de compromiso: una vez realizada, la empresa no puede retirar el capital y queda expuesta a que el regulador no reconozca su costo.}

\sol{84}{A}{Depreciación $=0.05\times1,000,000=50,000$.}

\sol{85}{C}{La parte del costo común asignada al agua es $0.60\times100,000=60,000$.}

\sol{86}{D}{Cargo fijo por usuario $=1,000,000/10,000=100$.}

\sol{87}{D}{Cargo volumétrico $=500,000/100,000=5$ por m$^3$.}

\sol{88}{A}{La fusión puede reducir el costo medio si las economías de escala son suficientes para compensar los costos de coordinación y gestión.}

\sol{89}{A}{Mejorar continuidad y presión exige inversión, mantenimiento y operación; la regulación debe reconocer esos costos para preservar la sostenibilidad.}

\sol{90}{D}{La conservación cuesta S/ 1.50 por m$^3$, menos que los S/ 2.50 de la nueva fuente; es la opción económicamente preferible antes de ampliar la oferta.}

\subsection*{Bloque D. Ampliación: Teoría y Aplicaciones Prácticas (91 a 180)}

\sol{91}{C}{La libre disposición implica que si $y \in Y$ y $y' \le y$, entonces $y' \in Y$: la firma siempre puede descartar excedentes de producto o utilizar más insumos para obtener la misma o menor producción.}
\sol{92}{A}{Si $0 \in Y$ e $y \in Y$, por convexidad $\lambda y + (1-\lambda)0 = \lambda y \in Y$ para todo $\lambda \in [0,1]$, lo cual define rendimientos no crecientes a escala (se puede reducir la escala sin pérdida de eficiencia).}
\sol{93}{D}{Al mantener $\bar{K}$ fijo, cada unidad adicional de $L$ dispone de menos capital por trabajador, por lo que desde cierto punto $\partial^2 f/\partial L^2 < 0$ ($PMg_L$ decreciente). Esto es un fenómeno de corto plazo distinto a los rendimientos a escala.}
\sol{94}{B}{La etapa I termina donde $PMe_L$ es máximo ($PMg_L = PMe_L$) y la etapa III comienza cuando $PMg_L < 0$. En la etapa II, $PMe_L$ cae y $0 \le PMg_L \le PMe_L$ (elasticidad $\varepsilon_{Q,L} \in [0,1]$).}
\sol{95}{D}{Por definición, $\varepsilon_{Q,L} = \frac{\partial Q}{\partial L}\frac{L}{Q} = \frac{PMg_L}{PMe_L}$. En el máximo de $PMe_L$ se cumple $PMg_L = PMe_L$, por lo que $\varepsilon_{Q,L} = 1$.}
\sol{96}{C}{Por el teorema de Euler con homogeneidad de grado 1, $Q = L\cdot PMg_L + K\cdot PMg_K \implies PMe_L = PMg_L + \frac{K}{L}PMg_K$. Si $PMg_L > PMe_L$, necesariamente $PMg_K < 0$.}
\sol{97}{A}{Derivando $PMe_L = Q/L$ respecto a $L$ se obtiene $\frac{d PMe_L}{d L} = \frac{1}{L}(PMg_L - PMe_L)$, cuyo signo es estrictamente negativo whenever $PMg_L < PMe_L$.}
\sol{98}{B}{La convexidad estricta de las isocuantas (cuasiconcavidad estricta de $f$) implica que, conforme aumenta $L$ y cae $K$ sobre la misma isocuanta, la pendiente en valor absoluto $TMST_{L,K} = PMg_L/PMg_K$ decrece.}
\sol{99}{C}{Por definición, $\sigma = \frac{d\ln(K/L)}{d\ln(TMST_{L,K})}$, midiendo la curvatura de la isocuanta y cuán fácilmente la firma sustituye factores cuando cambian sus precios relativos $w/r$.}
\sol{100}{D}{En Cobb-Douglas, $TMST_{L,K} = \frac{\beta}{\alpha}\frac{K}{L} \implies \ln(K/L) = \ln(\alpha/\beta) + \ln(TMST_{L,K})$, por lo que $\sigma = \frac{d\ln(K/L)}{d\ln(TMST_{L,K})} = 1$ independientemente del grado de homogeneidad $\alpha+\beta$.}
\sol{101}{B}{Al multiplicar $(K,L)$ por $t>0$, $Q(tK, tL) = t^\nu Q(K,L)$. Así, $\nu$ gobierna los rendimientos a escala ($\nu=1$ constantes, $\nu>1$ crecientes), mientras que $\rho$ determina $\sigma = 1/(1-\rho)$.}
\sol{102}{A}{Como los insumos son complementos perfectos ($\sigma = 0$), el costo mínimo siempre ocurre en el vértice de cada isocuanta ($aK = bL$), de modo que la razón óptima $K/L = b/a$ es invariante a $w/r$.}
\sol{103}{C}{Con isocuantas lineales de pendiente constante $TMST_{L,K} = b/a$, si $b/a > w/r$ solo se usa $L$; si $b/a < w/r$ solo se usa $K$. Solo cuando $b/a = w/r$ coinciden isocuanta e isocosto permitiendo combinaciones interiores.}
\sol{104}{B}{La Translog es una aproximación de Taylor de segundo orden en logaritmos que no restringe $\sigma = 1$ ni los rendimientos a escala a ser constantes, permitiendo capturar economías de escala y alcance variables.}
\sol{105}{D}{Por el teorema de Euler, $K\cdot PMg_K + L\cdot PMg_L = k Q$. Multiplicando por $p$, el gasto en factores $rK + wL = k(pQ) > pQ$ cuando $k > 1$, lo que implica pérdidas si se fija $p = CMg$ (problema clásico del monopolio natural).}
\sol{106}{A}{En $Q = K^{0.7}L^{0.7}$, $\partial^2 Q/\partial L^2 < 0$ y $\partial^2 Q/\partial K^2 < 0$ (rendimientos marginales decrecientes en corto plazo), pero $Q(tK, tL) = t^{1.4}Q(K,L)$ (rendimientos crecientes a escala en largo plazo).}
\sol{107}{B}{Evaluando $\frac{d\ln f(tK, tL)}{d\ln t}\big|_{t=1} = \frac{K f_K + L f_L}{Q} = \varepsilon_{Q,K} + \varepsilon_{Q,L}$, se obtiene la elasticidad de escala local.}
\sol{108}{C}{La transformación monótona $F$ no altera la $TMST_{L,K}$ (que solo depende de $K/L$ por ser $g$ homogénea), por lo que la razón óptima $K^*/L^*$ es constante a lo largo de la senda de expansión, aunque $F$ permita rendimientos a escala variables.}
\sol{109}{A}{Con rendimientos constantes a escala, $C(w,r,q) = c(w,r)q$. Si $p > c(w,r)$, $\pi = (p-c)q \to \infty$ (no hay máximo finito); si $p = c(w,r)$, $\pi = 0$ para todo $q \ge 0$; si $p < c(w,r)$, $q^* = 0$.}
\sol{110}{D}{Por una propiedad estándar de funciones homogéneas, si $f$ es homogénea de grado $k$, sus derivadas de primer orden son homogéneas de grado $k-1$. Si $k=1$, $PMg_K$ y $PMg_L$ son homogéneas de grado $0$ (solo dependen de $K/L$).}
\sol{111}{C}{Reordenando $\frac{PMg_L}{PMg_K} = \frac{w}{r}$ se obtiene $\frac{PMg_L}{w} = \frac{PMg_K}{r}$, es decir, la productividad marginal por sol gastado se iguala en todos los factores utilizados.}
\sol{112}{B}{Por el teorema de la envolvente, $\frac{\partial C(w,r,q)}{\partial q} = \frac{\partial \mathcal{L}}{\partial q} = \mu^*$, interpretándose $\mu^*$ como el costo marginal de aumentar la meta de producción en una unidad.}
\sol{113}{A}{Si la firma mantuviera fijos los insumos iniciales $(L^0, K^0)$, el costo crecería linealmente con $(w,r)$; al reoptimizar y sustituir hacia el factor más barato, $C(w,r,q) \le wL^0 + rK^0$, lo que implica concavidad en $(w,r)$.}
\sol{114}{D}{Duplicar todos los precios de insumos no cambia la pendiente de las isocostos ($w/r$) ni la isocuanta objetivo $q$, por lo que la cesta óptima de insumos no varía (grado $0$) y el costo total se duplica (grado $1$).}
\sol{115}{B}{En el largo plazo la firma minimiza costos sobre un conjunto más amplio (sin la restricción $K=\bar{K}$), por lo que $C_{LP} \le C_{CP}$; ambos coinciden exactamente en el nivel de $q$ para el cual $\bar{K}$ es el capital óptimo de largo plazo.}
\sol{116}{C}{Si escalar insumos por $t$ multiplica el producto por $t^k = q$, entonces se requiere escalar los insumos (y el costo) por $t = q^{1/k}$, resultando $C(w,r,q) = C(w,r,1) \cdot q^{1/k}$.}
\sol{117}{D}{Por el lema de Shephard, $L^c = \partial C/\partial w$, de modo que $\partial L^c/\partial w = \partial^2 C/\partial w^2 \le 0$ por concavidad de $C$ en precios (ley de la demanda condicional de factores).}
\sol{118}{A}{Para minimizar una función lineal ($wL+rK$) sujeta a $f(K,L)=q$, basta que el conjunto superior de la isocuanta sea convexo (cuasiconcavidad de $f$), sin requerir concavidad global de $f$.}
\sol{119}{B}{Por el teorema de la envolvente aplicado al problema de minimización de costos, $\frac{\partial C(w,r,q)}{\partial w} = L^c(w,r,q)$ y $\frac{\partial C(w,r,q)}{\partial r} = K^c(w,r,q)$.}
\sol{120}{C}{Dado que $L^c = \partial C/\partial w$ y $K^c = \partial C/\partial r$, las derivadas cruzadas coinciden con la segunda derivada mixta $\frac{\partial^2 C}{\partial w \partial r} = \frac{\partial^2 C}{\partial r \partial w}$.}
\sol{121}{D}{Multiplicar $(p,w,r)$ por $t>0$ no altera los planes óptimos y multiplica $\pi$ por $t$ (grado 1). Además, $\pi(p,w,r)$ es el supremo de funciones lineales en $(p,w,r)$, por lo que es necesariamente convexa en todos los precios.}
\sol{122}{A}{Aplicando el teorema de la envolvente a $\pi(p,w,r) = p q^* - w L^* - r K^*$ se obtiene directamente $\frac{\partial \pi}{\partial p} = q^*$ y $\frac{\partial \pi}{\partial w} = -L^*$.}
\sol{123}{B}{Por convexidad de $\pi$, sus segundas derivadas directas son no negativas: $\frac{\partial q^*}{\partial p} = \frac{\partial^2 \pi}{\partial p^2} \ge 0$ y $\frac{\partial L^*}{\partial w} = -\frac{\partial^2 \pi}{\partial w^2} \le 0$. No existe equivalente a un bien Giffen en demanda de factores.}
\sol{124}{C}{Las condiciones de segundo orden de $\max_{K,L} pf(K,L)-wL-rK$ requieren que el Hessiano de $f$ sea semidefinido negativo (concavidad local), lo que corresponde a rendimientos decrecientes a escala locales y $CMg$ creciente.}
\sol{125}{A}{Maximizar $\pi = pf(K,L)-wL-rK$ respecto a $L$ exige $p \cdot PMg_L = w$ (o $PMg_L = w/p$): la empresa contrata trabajo hasta que el valor monetario del producto adicional iguale su costo salarial.}
\sol{126}{D}{Al no existir restricción de presupuesto fijo, la función de beneficios es convexa en precios y $\frac{\partial L^*}{\partial w} = -\frac{\partial^2 \pi}{\partial w^2} \le 0$. Incluso si un factor es inferior, su demanda incondicional nunca sube al encarecerse.}
\sol{127}{C}{Como $\frac{d CMe}{dq} = \frac{1}{q}(CMg - CMe)$ y $\frac{d CVMe}{dq} = \frac{1}{q}(CMg - CVMe)$, los mínimos ocurren donde $CMg = CMe$ y $CMg = CVMe$ con $CMg$ creciente.}
\sol{128}{B}{Como $CMe = CVMe + CFMe$, en el mínimo de $CVMe$ se tiene $\frac{d CVMe}{dq} = 0$ pero $\frac{d CFMe}{dq} = -CF/q^2 < 0$, por lo que $\frac{d CMe}{dq} < 0$: el $CMe$ sigue disminuyendo y alcanza su mínimo a una escala mayor $q_2 > q_1$.}
\sol{129}{A}{Dado $CV = w L(q)$, derivando respecto a $q$ resulta $CMg = w \frac{dL}{dq} = \frac{w}{PMg_L}$. Análogamente, $CVMe = \frac{w}{PMe_L}$, de modo que el mínimo de $CVMe$ coincide con el máximo de $PMe_L$.}
\sol{130}{D}{Si cierra, su beneficio es $-CF$. Si produce $q>0$, su beneficio es $(p - CVMe)q - CF$. Por tanto, conviene producir siempre que $p \ge \min CVMe$, incluso con pérdidas, porque contribuye a cubrir parte de los costos fijos.}
\sol{131}{B}{Al cerrar se evita tanto $CV(q)$ como $F_c$ (perdiendo solo $CF_h$). Por tanto, operar es mejor que cerrar solo si $pq - CV(q) - F_c \ge 0$, es decir, $p \ge \min \frac{CV(q)+F_c}{q}$.}
\sol{132}{C}{La firma iguala $p = CMg$ (en su tramo creciente por condición de segundo orden) siempre que $p \ge \min CVMe$; para precios menores ofrece $q^* = 0$.}
\sol{133}{D}{Los activos específicos como tuberías subterráneas de una EPS tienen valor de reventa casi nulo fuera de su uso actual (costos hundidos), lo que crea barreras a la salida y riesgos de oportunismo regulatorio (hold-up).}
\sol{134}{A}{Como $CMeLP$ y $CMeCP$ son tangentes, tienen la misma pendiente en el punto de contacto. Si $CMeLP$ es decreciente, $CMeCP$ también debe tener pendiente negativa allí, es decir, opera a la izquierda de su mínimo.}
\sol{135}{C}{Como $C_{LP}(q) = C_{CP}(q; K^*(q))$, por el teorema de la envolvente $\frac{d C_{LP}}{dq} = \frac{\partial C_{CP}}{\partial q}\big|_{K=K^*(q)}$, es decir, $CMgLP$ y $CMgCP$ se cortan exactamente en la producción de tangencia $q_0$.}
\sol{136}{B}{Cuando $\varepsilon_{C,q} = CMgLP/CMeLP < 1$, un aumento del $1\%$ en el producto eleva los costos en menos del $1\%$, por lo que el costo medio de largo plazo desciende (economías de escala).}
\sol{137}{D}{La EME es el menor volumen de producción en el que la firma agota todas las economías de escala disponibles y alcanza el mínimo de su curva de costo medio de largo plazo ($CMeLP$).}
\sol{138}{A}{Al ser $C_{CP}(q) \ge C_{LP}(q)$ con igualdad en $q_0$, la función $C_{CP}(q) - C_{LP}(q)$ tiene un mínimo en $q_0$, por lo que su segunda derivada es no negativa: $\frac{\partial CMgCP}{\partial q} \ge \frac{d CMgLP}{dq}$ (principio de Le Chatelier).}
\sol{139}{B}{Porque en el largo plazo la firma puede ajustar todos sus factores (eliminando la restricción $K=\bar{K}$), su respuesta productiva ante un alza en $p$ es mayor o igual que en el corto plazo: $\frac{\partial q^*_{LP}}{\partial p} \ge \frac{\partial q^*_{CP}}{\partial p}$.}
\sol{140}{C}{Un insumo es inferior si al aumentar la escala de producción $q$ (manteniendo $w, r$ fijos), la firma sustituye hacia tecnologías más intensivas en otros factores y reduce su demanda condicional: $\partial L^c/\partial q < 0$.}
\sol{141}{A}{Por el teorema de Young sobre $C(w,r,q)$, $\frac{\partial CMg}{\partial w} = \frac{\partial^2 C}{\partial w \partial q} = \frac{\partial L^c}{\partial q}$. Si $L$ es inferior ($\partial L^c/\partial q < 0$), un alza en $w$ reduce el $CMg$.}
\sol{142}{D}{Al subir $w$: (1) el efecto sustitución reduce $L^c$ a $q$ fijo ($\le 0$); (2) para un insumo normal, el $CMg$ sube, la firma reduce $q^*$ y el efecto producto reduce aún más $L$. Así, $|\partial L^*/\partial w| \ge |\partial L^c/\partial w|$.}
\sol{143}{B}{Por homogeneidad de grado 0 de $L^c(w,r,q)$, el teorema de Euler implica $w\frac{\partial L^c}{\partial w} + r\frac{\partial L^c}{\partial r} = 0$. Como $\frac{\partial L^c}{\partial w} < 0$, necesariamente $\frac{\partial L^c}{\partial r} > 0$.}
\sol{144}{C}{La subaditividad exige $C(q) < \sum_{i=1}^k C(q_i)$ para toda partición no trivial $\sum q_i = q$. En el caso uniproducto, puede existir subaditividad incluso en un tramo donde $CMeLP$ ya empezó a crecer levemente.}
\sol{145}{D}{Las economías de alcance miden el ahorro de costos derivado de la producción conjunta respecto a la producción especializada separada: $S_C = \frac{C(q_1,0)+C(0,q_2)-C(q_1,q_2)}{C(q_1,q_2)} > 0$.}
\sol{146}{A}{En el caso multiproducto, las economías de escala rayo (o específicas) no bastan para la subaditividad: si existieran deseconomías de alcance, dos firmas especializadas podrían producir más barato que una firma integrada.}
\sol{147}{C}{Las economías de densidad reflejan que, dada una longitud de tuberías instalada, aumentar el volumen servido o el número de conexiones por kilómetro de red diluye los costos fijos de infraestructura y reduce el costo unitario.}
\sol{148}{B}{Si $CMe(q)$ es estrictamente decreciente hasta $q$, entonces $C(q) < \sum C(q_i)$ siempre (suficiente). Sin embargo, una firma puede ser subaditiva en un rango ligeramente posterior al mínimo de $CMe(q)$ (no necesaria).}
\sol{149}{A}{Con $CMg < CMe$, cobrar $p = CMg$ genera ingresos $p q = CMg \cdot q < CMe \cdot q = CT$, provocando pérdidas financieras. Por ello la regulación recurre a tarifas en dos partes, Ramsey-Boiteux o precios iguales a $CMe$.}
\sol{150}{D}{En el progreso neutral de Hicks, $Q = A(t)f(K,L)$: ambos productos marginales crecen en la misma proporción $A(t)$, por lo que $TMST_{L,K} = f_L/f_K$ no varía para una relación $K/L$ dada.}
\sol{151}{C}{El progreso neutral de Harrod es aumentador de trabajo ($F(K, A(t)L)$), mientras que el neutral de Solow es aumentador de capital ($F(A(t)K, L)$) y el de Hicks es $A(t)F(K,L)$.}
\sol{152}{B}{En Cobb-Douglas, $A K^\alpha L^{1-\alpha} = K^\alpha (A^{1/(1-\alpha)}L)^{1-\alpha} = (A^{1/\alpha}K)^\alpha L^{1-\alpha}$, por lo que un avance neutral en cualquiera de los tres sentidos puede reescribirse en los otros dos.}
\sol{153}{D}{Diferenciando logarítmicamente $Q = A f(K,L)$ y usando que en competencia $s_K = \varepsilon_{Q,K}$ y $s_L = \varepsilon_{Q,L}$, se obtiene $\hat{Q} = \hat{A} + s_K \hat{K} + s_L \hat{L}$, de donde $\hat{A} = \hat{Q} - s_K \hat{K} - s_L \hat{L}$.}
\sol{154}{A}{El producto medio es $PMe_L = 18L - L^2$. Derivando e igualando a cero: $18 - 2L = 0 \implies L = 9.0$, donde $PMe_L = PMg_L = 81$.}
\sol{155}{B}{De $PMg_L = 36L - 3L^2 = 0$ resulta $L = 12$. Evaluando en la función de producción: $Q(12) = 18(144) - 1728 = 2592 - 1728 = 864$ unidades.}
\sol{156}{C}{En $L=12$: $PMe_L = 20(12) - 12^2 = 96$ y $PMg_L = 40(12) - 3(12^2) = 48$. Luego $\varepsilon_{Q,L} = \frac{PMg_L}{PMe_L} = \frac{48}{96} = 0.50$.}
\sol{157}{A}{La tasa marginal de sustitución técnica es $TMST_{L,K} = \frac{PMg_L}{PMg_K} = \frac{0.75}{0.25}\frac{K}{L} = 3\left(\frac{16}{8}\right) = 6.0$.}
\sol{158}{D}{Con $\rho = -0.25$, la elasticidad de sustitución técnica es $\sigma = \frac{1}{1 - \rho} = \frac{1}{1 - (-0.25)} = \frac{1}{1.25} = 0.80$.}
\sol{159}{C}{Escalando por $t>0$: $Q(tK, tL) = 2[t^{0.5}(K^{0.5}+L^{0.5})]^3 = t^{1.5}Q(K,L)$. Al ser homogénea de grado $1.5 > 1$, exhibe rendimientos crecientes a escala.}
\sol{160}{B}{En $(2,3)$, $Q = 10$, $PMg_K = L + 2 = 5$ y $PMg_L = K = 2$. Luego $E = \frac{K\cdot PMg_K + L\cdot PMg_L}{Q} = \frac{2(5) + 3(2)}{10} = 1.60$.}
\sol{161}{A}{La tangencia exige $K/L = w/r = 4/16 \implies L = 4K$. En la isocuanta, $\sqrt{4K^2} = 2K = 20 \implies K^* = 10, L^* = 40$. Así, $C = 4(40) + 16(10) = S/\ 320.00$.}
\sol{162}{D}{En el vértice óptimo $2K = 5L = 40$, por lo que $K^* = 20$ y $L^* = 8$. El costo mínimo es $C = 10(8) + 6(20) = S/\ 200.00$.}
\sol{163}{B}{Como $PMg_L/w = 2/6 = 0.333 > PMg_K/r = 3/12 = 0.25$, conviene usar únicamente trabajo: $L^* = 60/2 = 30$, $K^* = 0$, con costo $6(30) = S/\ 180$.}
\sol{164}{C}{La tangencia implica $(K/L)^{0.5} = w/r = 1/4 \implies L = 16K$. Sustituyendo en $Q = 25$: $(K^{0.5} + 4K^{0.5})^2 = 25K = 25 \implies K^* = 1, L^* = 16$, con costo $1(16) + 4(1) = S/\ 20.00$.}
\sol{165}{A}{De $TMST_{L,K} = 2K/L = w/r = 2$ resulta $K = L$. En $Q = K^{0.5}L = L^{1.5} = 8$ se obtiene $L^c = K^c = 4$, y el costo es $8(4) + 4(4) = S/\ 48$.}
\sol{166}{D}{Por el lema de Shephard, $L^c = \frac{\partial C}{\partial w} = q (r/w)^{0.75} = 10(16/1)^{0.75} = 10(8) = 80.0$ unidades.}
\sol{167}{B}{Derivando respecto a $r$: $K^c = \frac{\partial C}{\partial r} = q^2\left(3 + \sqrt{w/r}\right) = 4(3 + 3/2) = 18.0$ unidades.}
\sol{168}{C}{El costo marginal es $CMg = \frac{\partial C}{\partial q} = 6q\sqrt{wr} = 6(5)\sqrt{36} = S/\ 180.00$.}
\sol{169}{A}{Igualando $p\cdot PMg_L = w$: $6\left(\frac{5}{\sqrt{L}}\right) = 15 \implies \sqrt{L} = 2 \implies L^* = 4$ y $Q^* = 20$. El beneficio es $\pi^* = 6(20) - 15(4) - 40 = S/\ 20$.}
\sol{170}{B}{Por simetría ($w=r=2$), $K=L$. La CPO $p\cdot PMg_L = w$ da $8L^{-0.5} = 2 \implies L^* = K^* = 16$, $Q^* = 4(16^{0.5}) = 16$ y $\pi^* = 8(16) - 2(16) - 2(16) = S/\ 64$.}
\sol{171}{D}{Por el lema de Hotelling, $q^*(p,w,r) = \frac{\partial \pi}{\partial p} = \frac{10p}{\sqrt{wr}} = \frac{10(18)}{\sqrt{36}} = 30.0$ unidades.}
\sol{172}{C}{Por el lema de Hotelling, $L^* = -\frac{\partial \pi}{\partial w} = \frac{2.5p^2}{w^{1.5}r^{0.5}} = \frac{2.5(144)}{8 \cdot 3} = 15.0$ unidades.}
\sol{173}{B}{El costo variable medio es $CVMe = 20 - 4q + 0.5q^2$. Minimizando: $-4 + q = 0 \implies q = 4$, donde $p_{\text{cierre}} = \min CVMe = 20 - 16 + 8 = S/\ 12.00$.}
\sol{174}{A}{Igualando $CMg = CMe$: $10 + 2q = \frac{100}{q} + 10 + q \implies q^2 = 100 \implies q = 10$, con $CMe_{\min} = 10 + 10 + 10 = S/\ 30$.}
\sol{175}{C}{El costo medio evitable es $\frac{CV(q)+F_c}{q} = \frac{16}{q} + 4 + q$, cuyo mínimo ocurre en $q = 4$ dando $p_{\min} = 4 + 4 + 4 = S/\ 12.00$.}
\sol{176}{D}{Minimizando respecto a $K$: $4 - q^2/K^2 = 0 \implies K^*(q) = q/2$. Sustituyendo, la envolvente es $C_{LP}(q) = 4q$, que en $q = 30$ vale $S/\ 120.00$.}
\sol{177}{B}{El costo medio es $CMeLP = q^2 - 12q + 60$. Igualando su derivada a cero: $2q - 12 = 0 \implies q = 6$, donde $CMeLP(6) = 36 - 72 + 60 = S/\ 24$.}
\sol{178}{C}{En $q = 10$: $C(10) = 250$, $CMe = 25$ y $CMg = 10 + 10 = 20$. Por tanto, $\varepsilon_{C,q} = \frac{CMg}{CMe} = \frac{20}{25} = 0.80 < 1$, indicando economías de escala.}
\sol{179}{D}{El costo conjunto es $C(30,30) = 400$, mientras que por separado $C(30,0) = 250$ y $C(0,30) = 250$ (suma $500$). Luego $S_C = \frac{500 - 400}{400} = 0.25$.}
\sol{180}{A}{Con un operador, $C(8) = 36 + 64 = S/\ 100$; con dos operadores simétricos, $2C(4) = 2(36 + 16) = S/\ 104$. Al ser $100 < 104$, la función sigue siendo subaditiva aun superando el mínimo de $CMe$.}
